\documentclass[a4paper]{article}

\usepackage[a4paper, top=8mm, bottom=8mm, left=8mm, right=8mm]{geometry}

\usepackage{polyglossia}
\setdefaultlanguage[babelshorthands=true]{russian}
\setotherlanguage{english}

\usepackage{fontspec}
\setmainfont{FreeSerif}
\newfontfamily{\russianfonttt}[Scale=0.7]{DejaVuSansMono}

\usepackage[tiny, compact]{titlesec}

\usepackage{titling}
\setlength{\droptitle}{-1cm}
\setlength{\parskip}{1.5ex}
\pretitle{\begin{center}\begin{bfseries}\Large}
\posttitle{\par\end{bfseries}\end{center}}
\preauthor{\begin{center}\normalsize}
\postauthor{\par\end{center}\vspace{-1.8cm}}

\usepackage{enumitem}
\setlist[enumerate]{topsep=0pt, itemsep=0.2pt, leftmargin=2em}
\setlist[itemize]{topsep=0pt, itemsep=0.2pt, leftmargin=2em}

\usepackage{hyperref}
\usepackage{bookmark}
\usepackage{csquotes}
\usepackage{amsmath}
\usepackage{graphicx}
\usepackage{float}


\title{Реализация свёртки изображений}

\author{Парфёнова Виктория Вадимовна}

\date{}

\begin{document}

\maketitle

\begin{flushright}
    Группа: \emph{25.Б42-мм}

    Кафедра: \emph{системного программирования СПбГУ}

    Научный руководитель: \emph{Пономарев Николай Алексеевич}
    
    Номер семестра практики: \emph{2}
\end{flushright} 

\section{Введение}

Свёртка используется для выделения признаков изображения (границ, текстуры), применения фильтров (размытие, повышение резкости, тиснение), уменьшения шума, а также в свёрточных нейронных сетях для распознавания объектов.

Алгоритм свёртки заключается в следующем: по изображению скользит ядро размером $k\times k$ с шагом 1. Ядро свёртки (фильтр)~--- это небольшая, обычно квадратная матрица коэффициентов, которая определяет эффект преобразования изображения. Для каждой позиции вычисляется сумма поэлементных произведений пикселей окна и коэффициентов ядра:
\[
        O[i,j] = \sum_{u=0}^{k-1} \sum_{v=0}^{k-1} I_{\text{pad}}[i+u, j+v] \cdot K[u,v],
\]
где:
$O[i,j]$~--- значение для пикселя нового изображения; $I_{\text{pad}}$~--- матрица размером с ядро. Центр матрицы~--- исходный пиксель; $K$~--- ядро.

Пример ядра (\texttt{sharpen\_3x3}~--- повышение резкости):
\[
\begin{pmatrix}
0 & -1 & 0 \\
-1 & 5 & -1 \\
0 & -1 & 0
\end{pmatrix}
\]

Обработка краёв необходима, потому что у граничных пикселей изображения нет полного набора соседей (например, у левого верхнего угла нет соседей слева и сверху). Если ничего не предпринимать, выходное изображение будет меньше исходного на размер ядра. Чтобы сохранить размерность, применяется дополнение (padding): добавляются искусственные пиксели (нули, отражённые, скопированные и т.п.) перед свёрткой.

\section{Цель работы}

Целью работы является изучение алгоритма двумерной свёртки, его программная реализация и сравнение производительности полученной реализации с реализацией из библиотеки OpenCV.

\section{Реализация}

Для реализации был выбран язык программирования Python 3.10, так как он предоставляет широкий набор библиотек (NumPy, Pillow, Matplotlib и др.), необходимых для реализации свёртки, визуализации результатов и сравнения производительности.

В рамках работы реализованы функции \texttt{convolution\_grayscale} и \texttt{convolution\_rgb}, которые принимают изображение (в виде массива NumPy), ядро свёртки и режим обработки краёв.

Алгоритм работы:
\begin{enumerate}
    \item Исходное изображение дополняется или обрезается согласно выбранному режиму обработки краёв.
    \item Выполняется свёртка (для цветных изображений операция выполняется независимо для каждого канала, после чего полученные каналы группируются обратно в трёхмерный массив).
    \item Функция возвращает итоговое изображение (в виде массива NumPy).
\end{enumerate}

Реализация решения осуществлялась с использованием также следующих библиотек:
\begin{enumerate}
    \item NumPy~--- для работы с операциями над двумерными массивами.
    \item Pillow~--- для загрузки и сохранения изображений, преобразования в массивы NumPy.
    \item Matplotlib~--- для визуализации результатов и построения графиков производительности.
    \item pytest~--- для тестирования.
    \item pytest-benchmark~--- для замера производительности и сравнения времени выполнения.
\end{enumerate}

Управление зависимостями выполнялось через инструмент uv\footnote{uv~--- URL: \url{https://docs.astral.sh/uv/} (дата обращения: 20.05.2026).}.

Доступны следующие ядра: 
\begin{enumerate}
    \item \texttt{sharpen\_3x3}~--- повышение резкости.
    \item \texttt{blur\_3x3}~--- усредняющее размытие.
    \item \texttt{gaussian\_blur\_3x3}~--- размытие по Гауссу.
    \item \texttt{highlighting\_vertical\_borders\_3x3}~--- выделение вертикальных границ.
    \item \texttt{highlighting\_horizontal\_borders\_3x3}~--- выделение горизонтальных границ.
    \item \texttt{embossing\_3x3}~--- эффект тиснения.
    \item \texttt{blur\_5x5}~--- усредняющее размытие (более плавное размытие, чем у $3\times 3$).
    \item \texttt{gaussian\_blur\_5x5}~--- размытие по Гауссу (более плавное размытие, чем у $3\times 3$).
\end{enumerate}

Доступны следующие варианты обработки края:
\begin{enumerate}
    \item \texttt{no\_padding}~--- игнорирование краёв (результат будет меньше исходного изображения на размер ядра).
    \item \texttt{zero\_padding}~--- дополнение нулями.
    \item \texttt{replicate\_padding}~--- копирование крайних пикселей.
    \item \texttt{mirror\_padding}~--- отражение без дублирования границы.
    \item \texttt{symmetric\_padding}~--- симметричное отражение с дублированием границы.
    \item \texttt{tile\_padding}~--- циклическое продолжение.
\end{enumerate}

\subsection{Запуск программы}

Выбор конкретного ядра, режима обработки краёв, входного изображения, выходной директории, количества каналов (чёрно-белое или цветное изображение) осуществляется через аргументы командной строки с использованием библиотеки Typer\footnote{Typer~--- URL: \url{https://typer.tiangolo.com/} (дата обращения: 20.05.2026).}. 

Пример запуска программы:
\begin{verbatim}
uv run -m src.main run --input-dir INPUT_DIR --output-dir OUTPUT_DIR --image-mode IMAGE_MODE --kernel KERNEL --padding PADDING
\end{verbatim}

Где:
\begin{itemize}
    \item \texttt{--input-dir}~--- путь к директории с входными изображениями;
    \item \texttt{--output-dir}~--- путь к директории для сохранения обработанных изображений;
    \item \texttt{--image-mode}~--- режим цветности: L (чёрно-белые изображения) или RGB (цветное изображение);
    \item \texttt{--kernel}~--- ядро свёртки (например, \texttt{blur\_3x3}, \texttt{sharpen\_3x3});
    \item \texttt{--padding}~--- режим обработки краёв (например, \texttt{zero\_padding}, \texttt{replicate\_padding}, \texttt{no\_padding}).
\end{itemize}

\subsection{Тестирование}

Для проверки корректности собственной реализации были написаны интеграционные тесты. Для создания эталонных изображений использовалась функция scipy.signal.convolve2d, так как именно она позволяет проверять режим \texttt{no\_padding}~--- обрезка краёв (в scipy.signal.convolve2d - это режим \texttt{valid})\footnote{SciPy API ~--- Signal processing (scipy.signal) ~--- convolve2d --- URL:\url{https://docs.scipy.org/doc/scipy/reference/generated/scipy.signal.convolve2d.html} (дата обращения: 21.05.2026).}. 
cv2.filter2D из OpenCV всегда возвращает изображение того же размера, что и входное\footnote{Function opencv::imgproc::filter\_2d~--- URL:\url{https://docs.opencv.org/4.x/dd/d6a/tutorial_js_filtering.html} (дата обращения: 21.05.2026).}, поэтому не позволяет проверить все варианты обработки краёв, реализованные в работе.

\section{Эксперименты}

Эксперименты проводились на системе со следующими характеристиками:
\begin{itemize}
    \item Процессор: Intel Core i5-1340P
    \item ОЗУ: 16 ГБ
    \item ОС: Ubuntu 22.04.5
\end{itemize}

Проведён эксперимент по сравнению скорости выполнения свёртки между учебной реализацией и библиотечной функцией OpenCV (cv2.filter2D).

Измерения проводились для квадратных изображений трёх размеров: $128\times 128$, $512\times 512$ и $2048\times 2048$ пикселей. Для каждого размера использовались следующие комбинации ядра свёртки и режима обработки краёв. 

Для цветных изображений: 
\begin{itemize}
    \item ядро \texttt{embossing\_3x3} с режимом \texttt{zero\_padding};
    \item ядро \texttt{gaussian\_blur\_5x5} с режимом \texttt{replicate\_padding};
    \item ядро \texttt{blur\_5x5} с режимом \texttt{mirror\_padding};
    \item ядро \texttt{embossing\_3x3} с режимом \texttt{tile\_padding}.
\end{itemize}
Для чёрно-белых изображений:
\begin{itemize}
    \item ядро \texttt{sharpen\_3x3} с режимом \texttt{no\_padding};
    \item ядро \texttt{blur\_3x3} с режимом \texttt{symmetric\_padding};
    \item ядро \texttt{gaussian\_blur\_3x3} с режимом \texttt{no\_padding};
    \item ядро \texttt{highlighting\_vertical\_borders\_3x3} с режимом \texttt{symmetric\_padding}.
\end{itemize}

\begin{figure}[H]
    \centering
    \includegraphics[width=0.9\textwidth]{images/benchmark\_results.pdf}
    \caption{Сравнение производительности библиотеки OpenCV и учебной реализации}
    \label{fig:benchmark_results}
\end{figure}

На графиках видно, что при росте размера изображения время работы учебной реализации увеличивается пропорционально размеру изображения ($O(height\times width)$). Функция cv2.filter2D из OpenCV демонстрирует аналогичную линейную зависимость от размера, но с гораздо меньшей константой (то есть на любом размере изображения OpenCV работает на порядки быстрее).

Различные режимы обработки краёв варьируют время вычислений в обоих случаях, так как, например, зеркальное отражение границ (\texttt{mirror\_padding}) требует сложных расчётов в отличие от простого заполнения нулями \texttt{zero\_padding}.

Обработка чёрно-белых изображений выполняется быстрее цветных для обоих алгоритмов. Это связано с сокращением числа каналов, что пропорционально уменьшает объём вычислений.
\begin{figure}[H]
    \centering
    \includegraphics[width=0.9\textwidth]{images/grayscale\_vs\_rgb.pdf}
    \caption{Сравнение производительности чёрно-белых и цветных изображений для библиотеки OpenCV и учебной реализации}
    \label{fig:grayscale_vs_rgb}
\end{figure}

Вывод, полученный из экспериментов: реализованная свертка уступает в производительности функции cv2.filter2D из OpenCV\footnote{Benchmark results description~--- URL: \url{https://github.com/roriess/image-convolution/blob/main/src/benchmark/benchmark_results_description.md} (дата обращения: 20.05.2026).}. Это связано с тем, что учебная свертка написана на языке Python с использованием вложенных циклов и NumPy, а pеализация OpenCV - на C++ и содержит низкоуровневые оптимизации, например, используется алгоритм с SIMD-инструкциями (команды, которые выполняют одну операцию сразу над несколькими элементами данных)\footnote{Function opencv::imgproc::filter\_2d~--- URL:\url{https://docs.opencv.org/4.x/dd/d6a/tutorial_js_filtering.html} (дата обращения: 21.05.2026).}.

\section{Заключение}

В ходе выполнения работы были достигнуты следующие цели:
\begin{enumerate}
    \item Изучена теория по свертке и необходимым для практической части пакетам.
    \item Реализована свертка для чёрно-белых и цветных изображений с возможным выбором ядра и варианта обработки края. 
    \item Поставлены эксперименты.
\end{enumerate}

Репозиторий с реализацией: \url{https://github.com/roriess/image-convolution}

\end{document}